\section{Preliminaries}
\label{sec:prelim}
%
We write $[n]\coloneqq \{1,\dots, n\}$ and for any set $S$ and integer $d$, let ${S \choose d}$ denote the set of $d$-element subsets of $S$.
%
For a finite constraint family $\confam$, we say its \emph{arity $r$} is the maximum arity of a function $f\in \confam$. Since in the constraints of $\SAT(\confam)$, we may use variables in arbitrary order, we use the following notation for convenience: For any $f:\{0,1\}^r \to \{0,1\}$ and partition $X_1,\dots, X_s$ of $[r]$, we write  
\[ f(\overbrace{\entry{x_1}}^{X_1},\dots, \overbrace{\entry{x_s}}^{X_s})\]
to denote the value of $f(u_1,\dots,u_r)$ where we plug in $x_j$ for each $u_i$ with $i\in X_j$. Correspondingly $g:\{0,1\}^d \to \{0,1\}$ can be obtained as a restriction of $f$ if and only if there is a partition $X_1,\dots, X_d, Z_0, Z_1$ of $[r]$ such that 
\[ g(x_1,\dots,x_d) = f(\overbrace{\entry{x_1}}^{X_1},\dots, \overbrace{\entry{x_d}}^{X_d}, \overbrace{\entry{0}}^{Z_0}, \overbrace{\entry{1}}^{Z_1}).\]
%
We say that an assignment $a: [n] \to \{0,1\}$ has \emph{weight} $k$ if $\sum_{i=1}^n a(i) = k$. Furthermore, we say that $a$ is dominated by an assignment $a': [n]\to \{0,1\}$, written $a\le a'$, if for all $i\in [n]$, we have $a(i) \le a'(i)$. For a subset $S\subseteq [n]$, we let $a_S$ denote the assignment that sets $a(i)=1$ if and only if $i\in S$. We let $\ones(a) \coloneqq \{ x_i \mid a(i) = 1\}$ denote the set of $1$-variables of $a$. For any constraint $C=f(\mathbf{x})$ where $\mathbf{x} = (x_{i_1},\dots,x_{i_r})$ with $i_1,\dots, i_r\in [n]$, we let $\vars(C) = \{x_{i_1},\dots,x_{i_r}\}$ denote the variable set involved in $C$.
%
%
All graphs considered in this paper are simple, i.e., we disallow multiple edges and self-loops. If $G=(V,E)$ is a directed graph, we call $S\subseteq V$ a \emph{closed} set if for all $(u,v)\in E$, we have that $u\in S$ implies that $v\in S$. We say that $v$ is a \emph{descendant} of $u$ if $v$ is reachable by a path from $u$ and let $D(u)$ denote the set of descendants of $u$ (including $u$ itself). Analogously, if $v$ is a descendant of $u$, we call $u$ an \emph{ascendant} of $v$. We extend the notation naturally to sets $S\subseteq V$ by defining $D(S) \coloneqq \bigcup_{u\in S} D(u)$.
For a graph $G=(V,E)$ with node weights $w: V\to \nats$ and $S\subseteq V$, we write $w(S) \coloneqq \sum_{v\in S}w(v)$. For any $S\subseteq V$, we let $G[S]$ denote the subgraph of $G$ induced by $S$, i.e., the subgraph obtained by deleting all vertices in $V\setminus S$ and adjacent edges.
%
\subsection{Hardness Assumptions}
%
Let $k$-clique denote the following problem: Given a (simple) undirected graph $G=(V,E)$, determine whether there is a \emph{clique} of size $k$, i.e., $S\subseteq V, |S|=k$ such that for all $\{u,v\} \in {S \choose 2}$ we have $\{u,v\} \in E$. A simple algorithm~\cite{NesetrilP85} solves $k$-clique in time $\Oh(n^{\omega/3 k})$ when $k$ is divisible by~$3$, which extends to time $\Oh(n^{\lfloor k/3\rfloor \omega + (k\bmod 3)})$ for arbitrary $k$ (for more precise bounds, see~\cite{EisenbrandG04}).
\label{sec:hardnessassumptions}
This running time is conjectured to be best possible, in the following sense.
%
\begin{hypo}[$k$-Clique Conjecture]
For no $c,\eps > 0$ and $f(k)$, there is an $f(k)n^{(\omega/3 - \epsilon)k+c}$-time algorithm for $k$-Clique.\footnote{Note: sometimes, the $k$-clique conjecture is stated as \[\inf \{ F \mid \text{$3k$-clique can be solved in time $n^{Fk+o(1)}$ for all (sufficiently large) constant $k$}\} = \omega,\] which can be seen to be equivalent to the above formulation via a standard self-reduction for $k$-clique.}
\end{hypo}
As without the use of matrix multiplication, no $\Oh(n^{(1-\eps)k+c})$-time algorithms are known, a variant of the conjecture postulates that there are even no $\Oh(n^{(1-\eps)k+c})$-time \emph{combinatorial} algorithms, i.e., algorithms avoiding the sophisticated algebraic techniques underlying current matrix multiplication algorithms. 
  
  By now, the $k$-clique conjecture has been used to explain hardness barriers in various contexts, such as the optimality of Valiant's parser for context-free grammar recognition~\cite{AbboudBVW18},  pattern matching in uncompressed and compressed strings~\cite{BringmannGL17, AbboudBBK17}, average-case hardness~\cite{Boix-AdseraBB19} and more. 
  For a more detailed discussion of this hardness assumption, we refer to~\cite{AbboudBVW18}.
%
  The $k$-clique problem naturally extends to hypergraphs: Given a $d$-uniform hypergraph $G=(V,E)$, the $d$-uniform $k$-HyperClique problem asks to determine whether there is a \emph{(hyper-)clique} of size $k$, i.e., $S\subseteq V, |S|=k$ such that for all subsets $S' \in {S \choose d}$, we have $S'\in E$. 
%
\begin{hypo}[$d$-Uniform $k$-HyperClique Conjecture]
Let $d\ge 3$. For no $c,\eps > 0$ and $f(k)$, there is an $f(k)n^{(1 - \epsilon)k+c}$-time algorithm for $d$-uniform $k$-HyperClique.
\end{hypo}
%
Similarly to the $k$-Clique conjecture, this hardness conjecture reveals hardness barriers in a number of contexts, such as hardness for problems on sparse graphs~\cite{LincolnVWW18}, for deciding or enumerating answers to first-order queries~\cite{BringmannFK19, CarmeliK19} and for the study of fine-grained average-case complexity~\cite{Boix-AdseraBB19}. It is known that it implies the Orthogonal Vectors conjecture~\cite{AbboudBDN18}, however, refuting this conjecture requires (at least) to give an $\Oh((2-\eps)^n)$-time exact algorithm for Max3SAT; for details and further discussion of the plausibility of this conjecture, we refer to \cite{LincolnVWW18}.   
%
